\section{Corpus jurídico delimitado para el cálculo de plazos}\label{sec:pruebas-corpus-juridico}

El corpus comprende investigación complementaria, su prórroga, dos modalidades de revocación y cuatro supuestos de apelación. Cada ficha distingue fuente normativa, interpretación adoptada, entradas ficticias y resultado esperado obtenido independientemente del motor. El caso de estudio corresponde a un órgano penal local de la Ciudad de México y utiliza notificación personal. Las fuentes fijan el fundamento y sus límites; las fechas de los ejemplos no se presentan como hechos de un expediente real ni requieren un precedente con las mismas fechas. El texto del CNPP empleado es la edición consolidada cuya última reforma indicada es del 28 de noviembre de 2025 \cite{cnpp_consolidado2025}. Su identificación no sustituye la comprobación de aplicación temporal a otro asunto.

\subsection{Calendario y canal de presentación}

El calendario seleccionado abarca del 1 de enero al 30 de noviembre de 2026 y se apoya en la Circular CJCDMX-21/2025, Acuerdo 76-29/2025, páginas 1 a 4 \cite{cjcdmx_calendario2026}. Excluye fines de semana e inhábiles publicados dentro de ese intervalo. El régimen penal conserva guardias y vacaciones escalonadas; por ello, la suspensión general de otras materias del 15 al 31 de julio no se traslada automáticamente a las UGJ. Un acuerdo particular no cubierto o una incidencia de aplicación desconocida impiden utilizar el resultado como vencimiento operativo. La ficha no presupone una revisión diaria de acuerdos posteriores.

Para los cinco recursos escritos se adopta la ventanilla ordinaria de atención al litigante. El artículo 4.I del Acuerdo 80-29/2025 establece atención de 09:00 a 15:00 de lunes a jueves y de 09:00 a 14:00 los viernes, salvo inhábiles \cite{cjcdmx_ugj2025}. El corte semanal usa esos cierres, desfase declarado de $-06{:}00$ y ausencia de cobertura ordinaria en sábado y domingo. Se excluyen sistemas del Ministerio Público, notificación electrónica y asuntos urgentes. El cierre representa disponibilidad del canal elegido, no una hora universal de caducidad o admisión por cualquier medio. El instante final de la audiencia oral y el cierre expreso de las resoluciones mensuales se conservan separados de este horario.

\subsection{Investigación complementaria y prórroga}

Los artículos 94 y 321 permiten una duración judicialmente determinada; el máximo es de dos meses cuando la pena máxima no excede dos años y de seis en caso contrario. El artículo 322 exige que la prórroga concedida y el período original respeten ese máximo \cite[arts.~94, 321--322]{cnpp_consolidado2025}. El amparo en revisión 630/2022 examina la imposibilidad de ampliar un máximo agotado, pero no establece una fórmula mensual universal \cite[párrs.~37--42 y 50--51]{scjn_ar630_2022}. Las fechas narradas en sus antecedentes tampoco justifican restar un día a toda suma mensual ni equiparar meses con treinta días.

Las dos fichas mensuales suponen una resolución ficticia que declara inicio, cantidad efectivamente concedida, cómputo de fecha a fecha por día homólogo y cierre a las 18:00 con desfase $-06{:}00$. No se proporciona una fecha final como entrada. Esta política es una premisa expresa del supuesto, no una regla mensual atribuida al texto general del CNPP. La ausencia de cantidad bloquea el cálculo aunque se conozca el máximo. En la prórroga también se identifican concesión, período original, ampliaciones previas y acumulado calificado; una petición aislada no equivale a concesión y la comparación no se presenta como suma automática del historial.

\begin{table}[htbp]
  \centering
  \small
  \caption{Entradas y resultados esperados de las fichas mensuales de 2026}
  \label{tab:corpus-mensual}
  \begin{tabularx}{\linewidth}{@{}L{2.5cm} X X@{}}
    \toprule
    Ficha & Entrada ficticia y condición & Resultado esperado \\
    \midrule
    Investigación & Inicio 12-01-2026; un mes concedido; máximo de dos meses. & Candidata 12-02-2026; cierre expreso equivalente a 13-02-2026 00:00 UTC. \\
    \addlinespace
    Prórroga & Mes inicial desde 12-01-2026; otro mes concedido desde 12-02-2026; acumulado calificado de dos meses. & Candidata 12-03-2026; cierre expreso equivalente a 13-03-2026 00:00 UTC. \\
    \bottomrule
  \end{tabularx}
\end{table}

Los contrastes de extremos mensuales utilizan un calendario sintético distinto, sin acreditar inhábiles oficiales futuros. Con la política expresa, 31-01-2027 más dos meses produce 31-03-2027 y 29-01-2028 más un mes produce 29-02-2028. El avance es directo al mes de destino, sin encadenar ajustes de febrero. En cambio, 31-01-2027 más un mes carece de homólogo y debe bloquearse si no existe una política expresa para resolverlo. La reserva afecta a ese negativo y a resoluciones insuficientes, no a todos los positivos mensuales. También deben bloquearse acumulados de tres meses bajo máximo de dos y de siete bajo máximo de seis, aunque la cantidad nueva aislada sea menor que el máximo. Unidades mezcladas o ajustes por inhábil que entren en conflicto con el límite acumulado quedan fuera de estas premisas.

\subsection{Revocación y apelaciones}

Las fichas escritas fijan notificación personal el miércoles 14 de octubre de 2026, efectos y primer día incluido el jueves 15. Para dos, tres y cinco días, la inclusión es la interpretación normativa expresamente seleccionada de los artículos 82, 84 y 94; las duraciones y los supuestos se identifican en los artículos 465, 466, 467 y 471 \cite{cnpp_consolidado2025}. La contradicción de tesis 308/2019 no unificó el inicio: declaró inexistente la contradicción al examinar porciones y supuestos distintos \cite[párrs.~23--35 y resolutivo único]{scjn_ct308_2019}. Por tanto, no se le atribuye una regla común para todas las apelaciones.

Para la sentencia de enjuiciamiento, el amparo directo en revisión 2742/2023 sitúa el inicio de los diez días en el día hábil posterior a la lectura y explicación, coincidente con efectos bajo las condiciones examinadas; exige disponibilidad de la versión escrita en esa audiencia \cite[párrs.~101--102, 116 y 146]{scjn_adr2742_2023}. Este criterio se utiliza en esa ficha sin extenderlo a las de dos, tres o cinco días. Las fechas siguientes se obtienen enumerando los días computables, con exclusión de los fines de semana del 17--18 y 24--25 de octubre.

\begin{table}[htbp]
  \centering
  \small
  \caption{Seis fichas de recursos y sus límites esperados en octubre de 2026}
  \label{tab:corpus-recursos}
  \begin{tabularx}{\linewidth}{@{}L{3.0cm} X L{3.0cm}@{}}
    \toprule
    Ficha & Supuesto controlado & Límite esperado \\
    \midrule
    Revocación oral & Trámite sin sustanciación; audiencia del 09-10, resolución 10:05:00 y cierre 10:50:00, con desfase $-06{:}00$. & Ancla 09-10, 10:50:00 local; 16:50:00 UTC. \\
    \addlinespace
    Revocación escrita & Resolución fuera de audiencia el 13-10, notificada en instalaciones el 14-10; dos días: 15 y 16. & 16-10, 14:00 local; 20:00 UTC. \\
    \addlinespace
    Apelación de auto de control & Negativa de anticipo de prueba; tres días: 15, 16 y 19. & 19-10, 15:00 local; 21:00 UTC. \\
    \addlinespace
    Apelación de sentencia de control & Sentencia definitiva en procedimiento abreviado; cinco días: 15, 16, 19, 20 y 21. & 21-10, 15:00 local; 21:00 UTC. \\
    \addlinespace
    Apelación sobre desistimiento de la acción penal & Resolución sobre desistimiento del MP ante tribunal de enjuiciamiento; tres días: 15, 16 y 19. & 19-10, 15:00 local; 21:00 UTC. \\
    \addlinespace
    Apelación de sentencia de enjuiciamiento & Lectura y explicación el 14-10, destinatario presente y texto escrito disponible; diez días: 15, 16, 19--23 y 26--28. & 28-10, 15:00 local; 21:00 UTC. \\
    \bottomrule
  \end{tabularx}
\end{table}

En la ficha oral, el contraste exige anterioridad estricta al cierre: 10:49:00 es anterior; 10:50:00 y 10:51:00 no lo son. La ausencia del cierre impide ese contraste, sin fabricar un plazo por días. Para las fichas escritas, incluir el día 14, desplazar el inicio al 16 o contar los fines de semana no reproduce los esperados. Asimismo, una sentencia no satisface la ficha de auto y el retiro de un recurso no constituye desistimiento de la acción penal. La distinción de supuestos precede al cómputo y este último no acredita por sí solo legitimación, procedencia o admisibilidad.

Un contraste adicional de revocación escrita comienza con notificación el 28 de octubre y efectos el 29. El primer día es jueves 29; se excluyen el viernes 30 publicado inhábil, sábado 31, domingo 1 de noviembre y lunes 2 publicado inhábil. El segundo día es martes 3 de noviembre, con cierre a las 15:00 local, equivalente a 21:00 UTC. Este caso distingue el calendario de la finalización semanal y detecta resultados incorrectos el 30 de octubre o el 2 de noviembre.

\subsection{Criterios de aceptación y alcance de la evidencia}

La aceptación técnica de cada positivo requiere observar la selección exacta de fuente, la cantidad y los días excluidos cuando correspondan, la candidata, la política de finalización y el instante UTC esperado. La captura debe conservar perfil, referencias y traza, así como calendario cuando proceda, mientras los negativos deben mostrar el motivo de bloqueo y ausencia de vencimiento operativo. La equivalencia entre una fecha civil y su representación UTC se contrasta separadamente para evitar confundir cambios de día con errores de cálculo. Los resultados manuales de las \cref{tab:corpus-mensual,tab:corpus-recursos} constituyen el oráculo del ensayo, no su salida precargada.

Las fichas delimitan entradas y observaciones exigibles; no registran por sí mismas pruebas ejecutadas, integración ni despliegue. Su reproducción documental no depende de una revisión externa obligatoria de los ocho casos. Las ejecuciones que las utilicen deben registrar separadamente entorno, resultados y límites, manteniendo las mediciones históricas de este capítulo. El fundamento y los supuestos completos se conservan en \rutarepo{docs/legal-corpus-investigation.md}, \rutarepo{docs/legal-corpus-resources.md} y \rutarepo{docs/legal-corpus-calendar.md}.

\subsection{Reproducción con servicios reales}

El 12 de octubre de 2026 UTC se ejecutó una campaña aislada con HTTP, PostgreSQL y Redis reales. Se conservaron ocho registros positivos, el contraste adicional de inhábiles y diez negativos. Los nueve instantes operativos coincidieron con los esperados y aparecieron en Agenda y alertas internas con la revisión y huella originales; los negativos conservaron motivo y ausencia de vencimiento. Reiniciar el servidor y restaurar la base en otra instancia mantuvo perfiles, fuentes, resultados, recibos e historial. La nueva autenticación mediante MFA recuperó acceso sin duplicar operaciones y la cadena de auditoría permaneció válida. Esta ejecución no utilizó datos productivos ni envió correo externo.

La aceptación de Qadra utilizó un servidor real y una cuenta aislada. Los recorridos de escritorio a 1440 píxeles y móvil a 390 píxeles contrastaron los ocho resultados, el caso de inhábiles, el calendario y las referencias capturados, los cierres semanales y la explicación por efecto de notificación ausente. Ambos aprobaron; el navegador informó 26.7 segundos. No se observaron errores JavaScript, escrituras del expediente durante las consultas ni desborde horizontal. La evidencia se conserva separada de la regresión integral y de los controles de integración y despliegue, que no se deducen de una aceptación local.
